\documentclass[12pt,a4paper]{article}
\usepackage{../latex/dsdlab-report}

\begin{document}
	\reporttitlepage{۴}{توصیف رفتاری}

	\section*{۴.۱ مقدمه و اهداف}
	هدف این آزمایش، آشنایی با توصیف رفتاری\LTRfootnote{Behavioral Modeling} در زبان توصیف سخت‌افزار \lr{Verilog} است. مطابق دستورکار، باید یک پشته\LTRfootnote{Stack} با عمق ۸ و پهنای ۴ بیت طراحی شود که ورودی‌های ساعت، بازنشانی، دادهٔ ورودی و دو فرمان درج و برداشت، و خروجی‌های دادهٔ خروجی و دو نشانک پر و خالی داشته باشد.

	پشته یک ساختار دادهٔ آخرین‌ورودی نخستین‌خروجی\LTRfootnote{Last In, First Out (LIFO)} است؛ یعنی هر برداشت، همان مقداری را بازمی‌گرداند که در آخرین درج نوشته شده است. برخلاف آزمایش سوم که استفاده از هر توصیفی جز جریان داده ممنوع بود، در این آزمایش دقیقاً عکس آن خواسته شده است: مدار باید با بلوک رویه‌ای \lr{always}، ساختار \lr{case} و انتساب غیرمسدودکننده\LTRfootnote{Non-Blocking Assignment} توصیف شود و رفتار مطلوب، نه ساختار دروازه‌ای آن، بیان گردد.

	\subsection*{روند انجام آزمایش}
	ابتدا معنای دقیق هر پایه و رفتار مورد انتظار در حالت‌های مرزی مشخّص و تصمیم‌های طراحی مستند شد. سپس پیمانهٔ پشته به‌صورت رفتاری و پارامتری نوشته شد و برای آن یک میز آزمون\LTRfootnote{Testbench} خودآزما\LTRfootnote{Self-Checking Testbench} ساخته شد که هفت سناریوی کارکردی و مرزی را می‌آزماید. اجرای شبیه‌سازی در \lr{ModelSim} با پرونده‌های فرمان خودکار انجام شد؛ یک بار با پیکربندی خواسته‌شدهٔ دستورکار و یک بار با پیکربندی دیگر برای نشان دادن پارامتری بودن واقعی طرح. در پایان، همان پروندهٔ سنتز‌پذیر در \lr{Quartus II} همگردانی شد و نمای سطح انتقال ثبات‌ها\LTRfootnote{Register Transfer Level (RTL) View}، نماهای فناوری\LTRfootnote{Technology Map Viewer}، گزارش مصرف منابع و تحلیل زمانی طرح بررسی شد.

	\section*{۴.۲ مشخّصات خواسته‌شده و تصمیم‌های طراحی}
	\subsection*{الف) واسط پیمانه}
	جدول زیر پایه‌های خواسته‌شده در دستورکار و نقش هر یک را گردآوری می‌کند.

	\begin{table}[H]
		\centering
		\small
		\begin{tabular}{|l|c|c|l|}
			\hline
			\textbf{نام} & \textbf{جهت} & \textbf{پهنا} & \textbf{نقش} \\ \hline
			\lr{Clk}      & ورودی & ۱ & ساعت مدار؛ همهٔ تغییرات در لبهٔ بالارونده \\ \hline
			\lr{RstN}     & ورودی & ۱ & بازنشانی ناهم‌زمان و فعّال‌صفر \\ \hline
			\lr{Data\_In} & ورودی & ۴ & داده‌ای که باید در پشته درج شود \\ \hline
			\lr{Push}     & ورودی & ۱ & فرمان درج \\ \hline
			\lr{Pop}      & ورودی & ۱ & فرمان برداشت \\ \hline
			\lr{Data\_Out}& خروجی & ۴ & داده‌ای که از پشته برداشته شده است \\ \hline
			\lr{Full}     & خروجی & ۱ & مقدار یک یعنی پشته پر است \\ \hline
			\lr{Empty}    & خروجی & ۱ & مقدار \textbf{صفر} یعنی پشته خالی است \\ \hline
		\end{tabular}
		\caption{پایه‌های پشته مطابق دستورکار}
	\end{table}

	\subsection*{ب) قرارداد نشانک خالی}
	یک نکتهٔ مهم که باید پیش از خواندن کد و شکل موج‌ها روشن شود، قطبیّت نشانک \lr{Empty} است. دستورکار به‌صراحت نوشته است \lr{Empty=0 indicates that the stack is empty}؛ یعنی این نشانک برخلاف انتظار متعارف، فعّال‌صفر است و هنگام خالی بودن پشته مقدار صفر و هنگام وجود دست‌کم یک عنصر مقدار یک می‌گیرد. در مقابل، نشانک \lr{Full} قطبیّت متعارف دارد و هنگام پر بودن پشته یک می‌شود.

	این قرارداد در سراسر طرح و میز آزمون رعایت شده است. بنابراین در شکل موج‌ها، بالا بودن \lr{Empty} به معنای «پشته خالی نیست» خوانده می‌شود و شرط \lr{if (Empty)} در کد، دقیقاً همان شرط «خالی نبودن پشته» است. برای پرهیز از ابهام، در متن این گزارش هرجا لازم بوده، معنای منطقی هر مقدار نیز آورده شده است.

	\subsection*{ج) حالت‌های تعریف‌نشده در دستورکار}
	دستورکار رفتار مدار را در چند وضعیّت مرزی مشخّص نکرده است. تصمیم‌های زیر گرفته و در کد و میز آزمون پیاده و آزموده شدند:
	\begin{itemize}
		\item \textbf{درج در پشتهٔ پر}: فرمان نادیده گرفته می‌شود. اشاره‌گر پشته\LTRfootnote{Stack Pointer} و محتوای حافظه دست‌نخورده می‌مانند تا داده‌های معتبر بازنویسی نشوند.
		\item \textbf{برداشت از پشتهٔ خالی}: فرمان نادیده گرفته می‌شود و \lr{Data\_Out} مقدار پیشین خود را نگه می‌دارد؛ زیرا هیچ دادهٔ معتبری برای بازگرداندن وجود ندارد و صفر کردن خروجی می‌توانست با یک دادهٔ واقعی صفر اشتباه گرفته شود.
		\item \textbf{درج و برداشت هم‌زمان}: به‌عنوان یک تصمیم طراحی، هیچ عملی انجام نمی‌شود. این انتخاب ساده‌ترین و قابل‌آزمون‌ترین رفتار ممکن است و از تعریف اولویّت دلبخواهی میان دو فرمان پرهیز می‌کند.
		\item \textbf{بازنشانی}: ناهم‌زمان و فعّال‌صفر است. تنها اشاره‌گر پشته و \lr{Data\_Out} صفر می‌شوند و آرایهٔ حافظه عمداً بازنشانی نمی‌شود؛ زیرا اشاره‌گر تضمین می‌کند هیچ‌گاه خانه‌ای خوانده نشود که پیش‌تر نوشته نشده باشد و بازنشانی ۳۲ فلیپ‌فلاپ حافظه تنها هزینهٔ اضافی می‌داشت.
	\end{itemize}

	چهار ترکیب ممکن دو فرمان و رفتار مدار در برابر هر یک، در جدول زیر خلاصه شده است.

	\begin{table}[H]
		\centering
		\small
		\begin{tabular}{|c|c|l|}
			\hline
			\textbf{\lr{Push}} & \textbf{\lr{Pop}} & \textbf{رفتار مدار} \\ \hline
			\(0\) & \(0\) & بدون فرمان؛ وضعیّت حفظ می‌شود \\ \hline
			\(1\) & \(0\) & درج در بالای پشته، تنها اگر پشته پر نباشد \\ \hline
			\(0\) & \(1\) & برداشت از بالای پشته، تنها اگر پشته خالی نباشد \\ \hline
			\(1\) & \(1\) & بی‌اثر؛ هیچ عملی انجام نمی‌شود \\ \hline
		\end{tabular}
		\caption{جدول تصمیم فرمان‌های درج و برداشت}
	\end{table}

	\section*{۴.۳ توصیف رفتاری پشته}
	\subsection*{الف) اشاره‌گر پشته و پهنای آن}
	هستهٔ طرح، ثبات \lr{Stack\_Pointer} است که تعداد عناصر معتبر پشته را نگه می‌دارد. با این تعریف، همین مقدار هم‌زمان نشانی نخستین خانهٔ آزاد نیز هست؛ ازاین‌رو درج در خانه‌ای از آرایهٔ \lr{Stack\_Mem} انجام می‌شود که نشانی آن برابر مقدار جاری اشاره‌گر است و سپس اشاره‌گر یک واحد افزایش می‌یابد؛ و برداشت از خانه‌ای انجام می‌شود که نشانی آن یک واحد کمتر از اشاره‌گر است و سپس اشاره‌گر یک واحد کاسته می‌شود. این انتخاب، منطق دو عمل را متقارن و ساده نگه می‌دارد.

	پهنای اشاره‌گر با رابطهٔ
	\[
	\text{\lr{PTR\_WIDTH}} = \bigl\lceil \log_2 (\text{\lr{DEPTH}}+1) \bigr\rceil
	\]
	محاسبه می‌شود و برای عمق ۸ برابر چهار بیت است. افزودن یک واحد به عمق در این رابطه ضروری است، زیرا اشاره‌گر باید بتواند \(\text{\lr{DEPTH}}+1\) مقدار متمایز، یعنی صفر تا ۸، را نمایش دهد. اگر به‌اشتباه از \(\lceil\log_2 8\rceil=3\) بیت استفاده می‌شد، مقدار ۸ که همان شرط پر بودن پشته است اصلاً نمایش‌پذیر نبود و شمارنده به‌جای رسیدن به آن، به صفر بازمی‌گشت.

	دو نشانک وضعیّت تنها از روی اشاره‌گر و به‌صورت ترکیبی ساخته می‌شوند:
	\[
	\text{\lr{Full}} = \bigl(\text{\lr{Stack\_Pointer}} = \text{\lr{DEPTH}}\bigr),
	\qquad
	\text{\lr{Empty}} = \bigl(\text{\lr{Stack\_Pointer}} \neq 0\bigr).
	\]
	وابسته بودن نشانک‌ها به یک منبع یگانه، امکان ناسازگاری میان وضعیّت واقعی پشته و نشانک‌های گزارش‌شده را از میان می‌برد.

	\subsection*{ب) بلوک رویه‌ای}
	کلّ رفتار مدار در یک بلوک \lr{always} حسّاس به لبهٔ بالاروندهٔ ساعت و لبهٔ پایین‌روندهٔ بازنشانی توصیف شده است. حضور \lr{negedge RstN} در فهرست حسّاسیّت همان چیزی است که بازنشانی را ناهم‌زمان\LTRfootnote{Asynchronous Reset} می‌کند. درون بلوک، ساختار \lr{case} روی الحاق دو بیت فرمان، چهار حالت جدول تصمیم را از هم جدا می‌کند و همهٔ انتساب‌ها از نوع غیرمسدودکننده‌اند تا رفتار ثبات‌های هم‌زمان به‌درستی مدل شود.

	طرح با دو پارامتر\LTRfootnote{Parameter} \lr{DEPTH} و \lr{WIDTH} نوشته شده است. مقادیر پیش‌فرض آن‌ها همان ۸ و ۴ خواسته‌شدهٔ دستورکار است، امّا هیچ جای کد به این دو عدد وابستگی سخت ندارد؛ این ادّعا در بخش شبیه‌سازی با اجرای پیکربندی دوم راستی‌آزمایی می‌شود.

	\vcode{code/stack.v}{توصیف رفتاری پشته}

	\section*{۴.۴ میز آزمون خودآزما}
	میز آزمون دو وظیفه\LTRfootnote{Task} کمکی دارد. وظیفهٔ \lr{push\_value} داده را در لبهٔ پایین‌رونده اعمال و فرمان درج را برای دقیقاً یک لبهٔ فعّال فعّال می‌کند؛ وظیفهٔ \lr{pop\_and\_check} یک برداشت انجام می‌دهد و مقدار بازگشتی را با مقدار انتظاری مقایسه و هر مغایرت را در شمارندهٔ خطا ثبت می‌کند. اعمال ورودی‌ها در لبهٔ پایین‌رونده و نمونه‌برداری با تأخیر یک نانوثانیه پس از لبهٔ بالارونده، از نقض شرط‌های برپایی و نگهداری و از رقابت خواندن و نوشتن در لحظهٔ لبه جلوگیری می‌کند.

	آزمون از هفت مرحلهٔ پیاپی تشکیل شده است:
	\begin{itemize}
		\item \textbf{مرحلهٔ ۱ --- بازنشانی}: پس از دو پالس ساعت با \lr{RstN} فعّال، بررسی می‌شود که پشته خالی باشد، نشانک پر فعّال نباشد و \lr{Data\_Out} صفر باشد.
		\item \textbf{مرحلهٔ ۲ --- برداشت از پشتهٔ خالی}: یک فرمان برداشت اعمال و بررسی می‌شود که وضعیّت پشته تغییر نکند.
		\item \textbf{مرحلهٔ ۳ --- پر کردن کامل}: دقیقاً \lr{DEPTH} مقدار پیاپی درج و پس از آن، فعّال شدن نشانک پر بررسی می‌شود.
		\item \textbf{مرحلهٔ ۴ --- درج در پشتهٔ پر}: یک درج دیگر اعمال و بررسی می‌شود که نشانک پر همچنان فعّال بماند و پشته سرریز نکند.
		\item \textbf{مرحلهٔ ۵ --- برداشت کامل و بررسی ترتیب}: همهٔ مقادیر برداشته و ترتیب آخرین‌ورودی نخستین‌خروجی آن‌ها یک‌به‌یک با مقدار انتظاری سنجیده می‌شود.
		\item \textbf{مرحلهٔ ۶ --- برداشت دوبارهٔ از پشتهٔ خالی}: افزون بر وضعیّت پشته، بررسی می‌شود که \lr{Data\_Out} نیز مقدار معتبر پیشین خود را نگه دارد.
		\item \textbf{مرحلهٔ ۷ --- درج و برداشت هم‌زمان}: پس از یک درج، هر دو فرمان با هم فعّال می‌شوند. در این مرحله، افزون بر خروجی‌ها، مقدار اشاره‌گر درونی نیز مستقیماً وارسی می‌شود.
	\end{itemize}

	وارسی مستقیم اشاره‌گر در مرحلهٔ هفتم تصمیم آگاهانه‌ای است. اگر تنها به خروجی‌ها بسنده می‌شد، مداری که به‌اشتباه هم درج و هم برداشت را انجام می‌دهد نیز از آزمون سربلند بیرون می‌آمد، چون اثر دو عمل بر اشاره‌گر یکدیگر را خنثی می‌کند. سنجش مستقیم \lr{Stack\_Pointer} این راه فرار را می‌بندد.

	\vcode{code/tb_stack.v}{میز آزمون خودآزمای پشته}

	\section*{۴.۵ اجرا در \lr{ModelSim}}
	برای آنکه اجرای شبیه‌سازی تکرارپذیر باشد و به دنبالهٔ فرمان‌های دستی وابسته نماند، سه پروندهٔ فرمان نوشته شد. پروندهٔ \lr{run.do} کتابخانهٔ پیشین را پاک و از نو می‌سازد، دو پروندهٔ توصیف سخت‌افزار را همگردانی می‌کند، شبیه‌ساز را روی میز آزمون بارگذاری و سپس پنجرهٔ موج را می‌سازد و شبیه‌سازی را تا پایان اجرا می‌کند. گزینهٔ \lr{-voptargs=+acc} از بهینه‌سازی سیگنال‌های داخلی جلوگیری می‌کند؛ بدون آن، اشاره‌گر و آرایهٔ حافظه در پنجرهٔ موج در دسترس نمی‌بودند و مرحلهٔ هفتم آزمون نیز نمی‌توانست اشاره‌گر را وارسی کند.

	\docode{code/run.do}{پرونده‌ی فرمان اجرای شبیه‌سازی با پیکربندی پیش‌فرض}

	پروندهٔ \lr{wave.do} سیگنال‌ها را در سه دستهٔ ورودی، خروجی و درونی به پنجرهٔ موج می‌افزاید و نمایش مقادیر عددی را به مبنای ده تغییر می‌دهد تا دنبال‌کردن مقادیر پشته آسان‌تر شود.

	\docode{code/wave.do}{پرونده‌ی پیکربندی پنجره‌ی موج}

	پروندهٔ سوم، \lr{run\_param.do}، همان طرح و همان میز آزمون را با پیکربندی دیگری اجرا می‌کند: عمق ۴ و پهنای ۸ بیت. مقادیر پارامترها با گزینهٔ \lr{-g} در زمان بارگذاری شبیه‌ساز جایگزین می‌شوند و هیچ تغییری در کد لازم نیست.

	\docode{code/run_param.do}{پرونده‌ی فرمان اجرای پیکربندی دوم برای راستی‌آزمایی پارامتری بودن طرح}

	تصویر زیر آغاز اجرای \lr{run.do} را نشان می‌دهد: ساخت کتابخانه، همگردانی موفّق دو پیمانه، بارگذاری شبیه‌ساز و افزودن سیگنال‌ها به پنجرهٔ موج.

	\shotw{screenshots/transcript_setup.png}{0.78}{آغاز اجرای پرونده‌ی فرمان و همگردانی موفّق پیمانه‌ها}

	خروجی کارکردی آزمون در ادامهٔ همان رونوشت چاپ می‌شود. هر سطر، فرمان انجام‌شده و مقدار دو نشانک را پس از آن گزارش می‌کند و در پایان پیام \lr{ALL TESTS PASSED} نشان می‌دهد که شمارندهٔ خطا صفر مانده است.

	\shotw{screenshots/transcript_result.png}{0.78}{خروجی کامل آزمون با پیکربندی عمق ۸ و پهنای ۴ بیت}

	دنبالهٔ چاپ‌شده دقیقاً همان است که انتظار می‌رفت: هشت درج پیاپی با مقادیر ۱ تا ۸، سپس یک درج نادیده‌گرفته‌شده در پشتهٔ پر، سپس هشت برداشت با ترتیب معکوس ۸ تا ۱، و سرانجام یک درج و یک برداشت پایانی برای آزمون فرمان هم‌زمان. سطر \lr{PUSH 8} نخستین سطری است که در آن \(\text{\lr{Full}}=1\) گزارش شده و سطر پایانی \lr{POP 1} نخستین سطری است که در آن نشانک خالی دوباره صفر، یعنی پشته خالی، شده است.

	\section*{۴.۶ تحلیل شکل موج‌ها}
	\subsection*{الف) سیگنال‌های افزوده‌شده به پنجره‌ی موج}
	جدول زیر ده سیگنال نمایش‌داده‌شده در پنجرهٔ موج و معنای آن‌ها را گردآوری می‌کند.

	\begin{table}[H]
		\centering
		\small
		\begin{tabular}{|l|l|l|}
			\hline
			\textbf{دسته} & \textbf{سیگنال} & \textbf{معنا} \\ \hline
			ورودی & \lr{Clk}            & ساعت با دورهٔ ۱۰ نانوثانیه \\ \hline
			ورودی & \lr{RstN}           & بازنشانی ناهم‌زمان فعّال‌صفر \\ \hline
			ورودی & \lr{Push}           & فرمان درج \\ \hline
			ورودی & \lr{Pop}            & فرمان برداشت \\ \hline
			ورودی & \lr{Data\_In}       & دادهٔ ورودی، در مبنای ده \\ \hline
			خروجی & \lr{Data\_Out}      & دادهٔ برداشته‌شده، در مبنای ده \\ \hline
			خروجی & \lr{Full}           & یک یعنی پشته پر است \\ \hline
			خروجی & \lr{Empty}          & صفر یعنی پشته خالی است \\ \hline
			درونی & \lr{Stack\_Pointer} & شمار عناصر معتبر و نشانی خانهٔ آزاد بعدی \\ \hline
			درونی & \lr{Stack\_Mem}     & آرایهٔ حافظهٔ پشته \\ \hline
		\end{tabular}
		\caption{سیگنال‌های پنجره‌ی موج و معنای آن‌ها}
	\end{table}

	\subsection*{ب) دنباله‌ی کامل رویدادها در پیکربندی خواسته‌شده}
	چون ساعت میز آزمون دورهٔ ۱۰ نانوثانیه دارد و نخستین لبهٔ بالارونده در لحظهٔ ۵ نانوثانیه رخ می‌دهد، لبه‌های فعّال در لحظه‌های \(5,15,25,\dots\) قرار می‌گیرند و می‌توان هر رویداد شکل موج را به یک لبهٔ مشخّص نسبت داد. جدول زیر کلّ شبیه‌سازی را لبه‌به‌لبه دنبال می‌کند؛ مقادیر ستون‌ها، وضعیّت پس از هر لبه‌اند.

	\begin{table}[H]
		\centering
		\footnotesize
		\begin{tabular}{|c|l|c|c|c|c|}
			\hline
			\textbf{زمان} & \textbf{فرمان} & \textbf{اشاره‌گر} & \textbf{\lr{Data\_Out}} & \textbf{\lr{Full}} & \textbf{\lr{Empty}} \\ \hline
			\(5,15\)   & بازنشانی فعّال      & \(0\) & \(0\)  & \(0\) & \(0\) \\ \hline
			\(25\)     & برداشت از خالی      & \(0\) & \(0\)  & \(0\) & \(0\) \\ \hline
			\(35\)     & درج \(1\)           & \(1\) & \(0\)  & \(0\) & \(1\) \\ \hline
			\(45\)     & درج \(2\)           & \(2\) & \(0\)  & \(0\) & \(1\) \\ \hline
			\(55\)     & درج \(3\)           & \(3\) & \(0\)  & \(0\) & \(1\) \\ \hline
			\(65\)     & درج \(4\)           & \(4\) & \(0\)  & \(0\) & \(1\) \\ \hline
			\(75\)     & درج \(5\)           & \(5\) & \(0\)  & \(0\) & \(1\) \\ \hline
			\(85\)     & درج \(6\)           & \(6\) & \(0\)  & \(0\) & \(1\) \\ \hline
			\(95\)     & درج \(7\)           & \(7\) & \(0\)  & \(0\) & \(1\) \\ \hline
			\(105\)    & درج \(8\)           & \(8\) & \(0\)  & \(1\) & \(1\) \\ \hline
			\(115\)    & درج \(15\) در پشتهٔ پر & \(8\) & \(0\)  & \(1\) & \(1\) \\ \hline
			\(125\)    & برداشت              & \(7\) & \(8\)  & \(0\) & \(1\) \\ \hline
			\(135\)    & برداشت              & \(6\) & \(7\)  & \(0\) & \(1\) \\ \hline
			\(145\)    & برداشت              & \(5\) & \(6\)  & \(0\) & \(1\) \\ \hline
			\(155\)    & برداشت              & \(4\) & \(5\)  & \(0\) & \(1\) \\ \hline
			\(165\)    & برداشت              & \(3\) & \(4\)  & \(0\) & \(1\) \\ \hline
			\(175\)    & برداشت              & \(2\) & \(3\)  & \(0\) & \(1\) \\ \hline
			\(185\)    & برداشت              & \(1\) & \(2\)  & \(0\) & \(1\) \\ \hline
			\(195\)    & برداشت              & \(0\) & \(1\)  & \(0\) & \(0\) \\ \hline
			\(205\)    & برداشت از خالی      & \(0\) & \(1\)  & \(0\) & \(0\) \\ \hline
			\(215\)    & درج \(15\)          & \(1\) & \(1\)  & \(0\) & \(1\) \\ \hline
			\(225\)    & درج و برداشت هم‌زمان & \(1\) & \(1\)  & \(0\) & \(1\) \\ \hline
			\(235\)    & برداشت              & \(0\) & \(15\) & \(0\) & \(0\) \\ \hline
		\end{tabular}
		\caption{دنباله‌ی کامل رویدادهای شبیه‌سازی با عمق ۸ و پهنای ۴ بیت؛ زمان‌ها بر حسب نانوثانیه}
	\end{table}

	\shot{screenshots/wave_default.png}{شکل موج کامل شبیه‌سازی با عمق ۸ و پهنای ۴ بیت}

	\subsection*{ج) تفسیر تک‌تک ویژگی‌های شکل موج}
	با کنار هم گذاشتن جدول بالا و شکل موج، همهٔ ویژگی‌های دیدنی در نما توضیح‌پذیر می‌شوند:

	\begin{itemize}
		\item \textbf{بازهٔ بازنشانی}: تا لحظهٔ ۲۰ نانوثانیه \lr{RstN} پایین است و اشاره‌گر و \lr{Data\_Out} روی صفر می‌مانند. چون بازنشانی ناهم‌زمان است، اشاره‌گر حتّی پیش از نخستین لبهٔ فعّال نیز مقدار معیّن صفر دارد و هیچ بازهٔ نامعیّنی در آن دیده نمی‌شود.
		\item \textbf{مقدار نامعیّن آرایهٔ حافظه}: ردیف \lr{Stack\_Mem} در آغاز شبیه‌سازی مقدار \lr{xxxx} دارد و تنها خانه‌هایی که نوشته می‌شوند به‌تدریج معیّن می‌گردند. این دقیقاً همان تصمیم طراحی است که در بخش ۴.۲ توضیح داده شد: آرایهٔ حافظه بازنشانی نمی‌شود، امّا چون اشاره‌گر هرگز به خانهٔ نانوشته اجازهٔ خواندن نمی‌دهد، این مقدار نامعیّن هیچ‌گاه به خروجی راه نمی‌یابد.
		\item \textbf{پلّهٔ صعودی اشاره‌گر}: میان لحظه‌های ۳۵ تا ۱۰۵، اشاره‌گر در هشت پلّهٔ یکسان از صفر به ۸ می‌رسد و ردیف \lr{Stack\_Mem} به‌ترتیب مقادیر \(0001\) تا \(1000\) را در خود جای می‌دهد.
		\item \textbf{نشانک خالی}: \lr{Empty} دقیقاً هم‌زمان با نخستین درج بالا می‌رود و تا لحظهٔ ۱۹۵ نانوثانیه که اشاره‌گر دوباره صفر می‌شود بالا می‌ماند. با توجّه به قرارداد فعّال‌صفر دستورکار، این یعنی «پشته از نخستین درج تا آخرین برداشت خالی نبوده است».
		\item \textbf{نشانک پر}: \lr{Full} تنها در بازهٔ ۱۰۵ تا ۱۲۵ نانوثانیه بالاست؛ یعنی از لحظه‌ای که هشتمین عنصر درج می‌شود تا لحظه‌ای که نخستین برداشت انجام می‌گیرد.
		\item \textbf{درج نادیده‌گرفته‌شده}: در لحظهٔ ۱۱۵ نانوثانیه \lr{Push} بالاست و \lr{Data\_In} برابر ۱۵ است، امّا اشاره‌گر روی ۸ ثابت می‌ماند و \lr{Stack\_Mem} تغییری نمی‌کند. این بخش، رفتار درست مدار در برابر سرریز\LTRfootnote{Overflow} را نشان می‌دهد.
		\item \textbf{پلّهٔ نزولی و ترتیب معکوس}: میان لحظه‌های ۱۲۵ تا ۱۹۵، اشاره‌گر در هشت پلّه به صفر بازمی‌گردد و \lr{Data\_Out} به‌ترتیب مقادیر ۸، ۷، ۶، ۵، ۴، ۳، ۲ و ۱ را نشان می‌دهد. این معکوس بودن ترتیب، همان ویژگی آخرین‌ورودی نخستین‌خروجی پشته است.
		\item \textbf{برداشت نادیده‌گرفته‌شده}: در لحظهٔ ۲۰۵ نانوثانیه \lr{Pop} بالاست، امّا اشاره‌گر روی صفر می‌ماند و \lr{Data\_Out} همچنان مقدار ۱ را نگه می‌دارد؛ یعنی مدار در برابر تهی‌روی\LTRfootnote{Underflow} نیز پایدار است و خروجی خود را با داده‌ای نامعتبر آلوده نمی‌کند.
		\item \textbf{فرمان هم‌زمان}: در لحظهٔ ۲۲۵ نانوثانیه هر دو سیگنال \lr{Push} و \lr{Pop} با هم بالا هستند. اشاره‌گر روی ۱ ثابت می‌ماند و \lr{Data\_Out} تغییر نمی‌کند؛ سپس برداشت پایانی در لحظهٔ ۲۳۵ مقدار ۱۵ را بازمی‌گرداند، که ثابت می‌کند مقدار درج‌شده پیش از فرمان هم‌زمان دست‌نخورده باقی مانده بود.
		\item \textbf{تأخیر یک‌پالسی نشانک‌ها}: نشانک‌ها ترکیبی‌اند و از اشاره‌گر ساخته می‌شوند؛ ازاین‌رو وضعیّت پس از آخرین عمل انجام‌شده را نشان می‌دهند، نه وضعیّت میان لبه را. برای نمونه، \lr{Full} در همان لبه‌ای بالا می‌رود که هشتمین درج در آن کامل می‌شود.
	\end{itemize}

	\subsection*{د) راستی‌آزمایی پارامتری بودن طرح}
	برای نشان دادن اینکه طرح به مقادیر ۸ و ۴ وابستگی سخت ندارد، همان کد و همان میز آزمون با عمق ۴ و پهنای ۸ بیت اجرا شد. رونوشت این اجرا نشان می‌دهد که مقادیر پارامترها در زمان بارگذاری جایگزین شده‌اند.

	\shotw{screenshots/transcript_param_setup.png}{0.78}{بارگذاری شبیه‌ساز با پارامترهای عمق ۴ و پهنای ۸ بیت}

	\shotw{screenshots/transcript_param_result.png}{0.78}{خروجی کامل آزمون با پیکربندی عمق ۴ و پهنای ۸ بیت}

	در این اجرا، پشته پس از چهار درج پر می‌شود، مقدار مرزی \(255\) که بزرگ‌ترین عدد نمایش‌پذیر در هشت بیت است به‌جای \(15\) به کار می‌رود، و کلّ آزمون در ۱۵۶ نانوثانیه به‌جای ۲۳۶ نانوثانیه پایان می‌یابد. هر هفت مرحلهٔ آزمون بدون تغییر در کد میز آزمون و بدون هیچ خطا اجرا شدند؛ این یعنی پارامتری بودن طرح صرفاً ادّعا نیست و در عمل نیز کار می‌کند. شکل موج این اجرا نیز همان الگوی پیشین را با مقیاس کوچک‌تر نشان می‌دهد.

	\shot{screenshots/wave_param.png}{شکل موج شبیه‌سازی با عمق ۴ و پهنای ۸ بیت}

	\section*{۴.۷ سنتز و پیاده‌سازی نرم‌افزاری}
	\subsection*{الف) همگردانی و مصرف منابع}
	تنها پروندهٔ سنتز‌پذیر طرح، یعنی \lr{stack.v}، در \lr{Quartus II 13.0.1} و برای خانوادهٔ \lr{Cyclone IV GX} همگردانی شد. انتخاب قطعه روی حالت خودکار گذاشته شد و ابزار قطعهٔ \lr{EP4CGX15BF14C6} را برگزید. همگردانی از مسیر \lr{Processing > Start} آغاز شد.

	\shotw{screenshots/quartus_menu.png}{0.72}{آغاز فرایند همگردانی و سنتز در \lr{Quartus II}}

	خلاصهٔ گزارش پس از مرحلهٔ تحلیل و سنتز، برآورد اوّلیّهٔ ۷۷ عنصر منطقی را نشان می‌دهد. در این مرحله هنوز جانمایی انجام نشده و قطعهٔ هدف نیز نهایی نشده است؛ ازاین‌رو درصد مصرف منابع گزارش نمی‌شود.

	\shotw{screenshots/flow_summary_synthesis.png}{0.9}{خلاصه‌ی گزارش پس از تحلیل و سنتز}

	پس از همگردانی کامل، جانمایی و مسیریابی نیز انجام می‌شود و ابزار می‌تواند بخشی از توابع ترکیبی و ثبات‌ها را در یک عنصر منطقی مشترک بسته‌بندی کند. به همین دلیل شمار عناصر منطقی در گزارش نهایی کمتر از برآورد مرحلهٔ سنتز است.

	\shotw{screenshots/flow_summary_compile.png}{0.9}{خلاصه‌ی گزارش پس از نخستین همگردانی کامل}

	\shotw{screenshots/flow_summary_final.png}{0.9}{خلاصه‌ی گزارش در آخرین همگردانی کامل؛ مرجع اعداد این گزارش}

	مرجع اعداد این گزارش، آخرین همگردانی کامل است: ۶۱ عنصر منطقی شامل ۴۵ تابع ترکیبی و ۴۰ ثبات اختصاصی، و ۱۴ پایه. مصرف عناصر منطقی کمتر از یک درصد ظرفیّت قطعه و مصرف پایه‌ها ۱۷ درصد است. اختلاف اندک شمار عناصر منطقی میان دو اجرای کامل، از تفاوت بسته‌بندی و جانمایی در هر اجرا ناشی می‌شود و ماهیّت طرح را تغییر نمی‌دهد.

	دو عدد این گزارش را می‌توان مستقیماً از روی کد راستی‌آزمایی کرد. شمار ثبات‌ها باید برابر باشد با هشت خانهٔ چهاربیتی حافظه، چهار بیت اشاره‌گر و چهار بیت \lr{Data\_Out}:
	\[
	8\times4 + 4 + 4 = 40,
	\]
	که دقیقاً با عدد گزارش‌شده می‌خواند. شمار پایه‌ها نیز برابر است با هشت پایهٔ ورودی، یعنی ساعت، بازنشانی، چهار بیت دادهٔ ورودی و دو فرمان، به‌علاوهٔ شش پایهٔ خروجی، یعنی چهار بیت دادهٔ خروجی و دو نشانک؛ در مجموع ۱۴ پایه.

	\subsection*{ب) نماهای سطح انتقال ثبات‌ها و فناوری}
	نمای سطح انتقال ثبات‌ها ساختار مورد انتظار را نشان می‌دهد: آرایهٔ حافظه به‌صورت یک بلوک با برچسب \lr{SYNC\_RAM}، ثبات اشاره‌گر به‌همراه دو جمع‌کننده برای افزایش و کاهش آن، مقایسه‌کننده‌های تولید نشانک‌ها و ثبات خروجی \lr{Data\_Out}. مسیر فرمان‌ها از پایه‌های \lr{Push} و \lr{Pop} به تسهیم‌کننده‌های انتخاب مقدار بعدی اشاره‌گر نیز به‌روشنی دیده می‌شود.

	\shot{screenshots/rtl_viewer.png}{نمای سطح انتقال ثبات‌های پشته}

	نکتهٔ درخور توجّه آن است که هرچند آرایهٔ حافظه در این نما به‌صورت حافظهٔ هم‌گام\LTRfootnote{Synchronous RAM} استنتاج شده است، گزارش منابع مقدار بیت‌های حافظهٔ مصرف‌شده را صفر اعلام می‌کند. دلیل آن است که کلّ حافظهٔ پشته تنها \(8\times4=32\) بیت است و ابزار ترجیح داده به‌جای اشغال یک بلوک حافظهٔ اختصاصی، آن را با ثبات‌های معمولی پیاده کند؛ همان ۴۰ ثباتی که در گزارش منابع دیده می‌شود.

	نماهای فناوری پس از نگاشت و پس از جانمایی نیز نشان می‌دهند که این ساختار به جدول‌های جست‌وجو و فلیپ‌فلاپ‌های واقعی قطعه نگاشت و مسیریابی شده است.

	\shot{screenshots/post_mapping.png}{نمای \lr{Technology Map Viewer (Post-Mapping)} پشته}
	\shot{screenshots/post_fitting.png}{نمای \lr{Technology Map Viewer (Post-Fitting)} پشته}

	\subsection*{ج) تحلیل زمانی}
	برخلاف برخی آزمایش‌های پیشین، این بار برای طرح یک پروندهٔ قیود طراحی\LTRfootnote{Synopsys Design Constraints} نوشته شد تا تحلیل زمانی معنادار باشد. این پرونده یک ساعت با دورهٔ ۲۰ نانوثانیه، معادل ۵۰ مگاهرتز، روی پایهٔ \lr{Clk} تعریف می‌کند.

	\docode{code/stack.sdc}{پرونده‌ی قیود طراحی و تعریف ساعت}

	تحلیل‌گر زمانی\LTRfootnote{TimeQuest Timing Analyzer} با مدل کند در دمای ۸۵ درجه و ولتاژ ۱۲۰۰ میلی‌ولت، یعنی بدترین شرایط کاری، اجرا شد.

	\shotw{screenshots/timequest_summary.png}{0.9}{خلاصه‌ی تحلیل‌گر زمانی و مدل تأخیر به‌کاررفته}

	خلاصهٔ برپایی نشان می‌دهد که حاشیهٔ زمانی\LTRfootnote{Slack} برای ساعت \lr{Clk} برابر \(16.037\) نانوثانیه و مجموع حاشیه‌های منفی\LTRfootnote{Total Negative Slack} برابر صفر است. مثبت بودن حاشیه و صفر بودن مجموع حاشیه‌های منفی یعنی طرح قید ۵۰ مگاهرتزی را با اطمینان برآورده می‌کند و برخلاف حالتی که هیچ قیدی تعریف نشده باشد، هیچ هشدار زمانی بی‌مبنایی تولید نمی‌شود.

	\shotw{screenshots/setup_summary.png}{0.8}{خلاصه‌ی تحلیل برپایی و حاشیه‌ی زمانی مثبت}

	از همین اعداد می‌توان بیشینه‌بسامد\LTRfootnote{Maximum Frequency (Fmax)} را نیز بازسازی کرد: بلندترین مسیر داده برابر \(20-16.037=3.963\) نانوثانیه است و وارون آن حدود ۲۵۲ مگاهرتز می‌شود. گزارش رسمی نیز مقدار \(252.33\,\mathrm{MHz}\) را نشان می‌دهد.

	\shotw{screenshots/fmax_summary.png}{0.9}{گزارش بیشینه‌بسامد و بیشینه‌بسامد مقیّد}

	با این حال، بیشینه‌بسامد مقیّد\LTRfootnote{Restricted Fmax} برابر \(250\,\mathrm{MHz}\) گزارش شده است. یادداشت خودِ گزارش علّت را روشن می‌کند: این حدّ نه از منطق درونی طرح، بلکه از محدودیّت نرخ تغییر پایه‌های ورودی و خروجی\LTRfootnote{Maximum I/O Toggle Rate} قطعه ناشی می‌شود. بنابراین بسامد کاری قابل‌اتّکا برای این پیاده‌سازی حدّاکثر ۲۵۰ مگاهرتز است و هم قید ۵۰ مگاهرتزی پروندهٔ قیود و هم ساعت ۱۰۰ مگاهرتزی میز آزمون، حاشیهٔ فراوان دارند.

	با توجّه به برگزاری مجازی آزمایشگاه، ارزیابی این آزمایش بر پایهٔ شبیه‌سازی، سنتز و تحلیل زمانی نرم‌افزاری انجام شده و در پرونده‌های پیکربندی پروژه تخصیص پایه‌ای برای برد ثبت نشده است؛ ازاین‌رو ادّعایی دربارهٔ آزمون فیزیکی روی برد مطرح نمی‌شود.

	\section*{۴.۸ نتایج}
	\begin{itemize}
		\item پشته‌ای با عمق ۸ و پهنای ۴ بیت و با همان پایه‌های خواسته‌شدهٔ دستورکار، تنها با توصیف رفتاری و در یک بلوک \lr{always} پیاده‌سازی شد.
		\item قرارداد فعّال‌صفر بودن نشانک \lr{Empty}، همان‌گونه که دستورکار تصریح کرده بود، در طرح و در میز آزمون رعایت شد.
		\item هر چهار ترکیب فرمان‌های درج و برداشت، و هر دو حالت مرزی درج در پشتهٔ پر و برداشت از پشتهٔ خالی، آزموده شدند و مدار در همهٔ آن‌ها وضعیّت خود را سالم نگه داشت.
		\item ترتیب آخرین‌ورودی نخستین‌خروجی برای هر هشت مقدار درج‌شده، یک‌به‌یک و به‌صورت خودکار وارسی شد.
		\item آزمون فرمان هم‌زمان، افزون بر خروجی‌ها، مقدار اشاره‌گر درونی را نیز وارسی کرد تا نتیجه به‌اشتباه قبول اعلام نشود.
		\item همان کد و همان میز آزمون با عمق ۴ و پهنای ۸ بیت نیز بدون هیچ تغییری و بدون خطا اجرا شد و پارامتری بودن طرح را تأیید کرد.
		\item سنتز با ۶۱ عنصر منطقی، ۴۰ ثبات و ۱۴ پایه انجام شد؛ شمار ثبات‌ها و پایه‌ها دقیقاً با محاسبهٔ نظری از روی کد می‌خواند.
		\item تحلیل زمانی با قید ۵۰ مگاهرتزی حاشیهٔ مثبت \(16.037\) نانوثانیه داد و بیشینه‌بسامد \(252.33\,\mathrm{MHz}\) با حدّ مقیّد \(250\,\mathrm{MHz}\) گزارش شد.
	\end{itemize}

	\section*{نتیجه‌گیری}
	در این آزمایش یک پشتهٔ سخت‌افزاری با توصیف رفتاری ساخته شد. مهم‌ترین درس آن، فاصلهٔ میان «بیان رفتار» و «بیان ساختار» بود: کلّ منطق پشته در چند ده سطر و بدون هیچ اشاره‌ای به دروازه یا فلیپ‌فلاپ نوشته شد، و سنتز‌کننده خود آرایهٔ حافظه، شمارنده و مقایسه‌کننده‌های لازم را استنتاج کرد. در کنار آن، دقّت در جزئیّاتی که دستورکار تعیین نکرده بود — قطبیّت نشانک خالی، رفتار در برابر سرریز و تهی‌روی، و معنای فرمان هم‌زمان — تعیین‌کنندهٔ درستی طرح بود؛ چنان‌که پهنای اشاره‌گر نیز اگر یک بیت کمتر گرفته می‌شد، شرط پر بودن پشته اصلاً نمایش‌پذیر نمی‌بود. میز آزمون خودآزما با پوشش هفت سناریو، اجرای دوبارهٔ همان آزمون با پارامترهای متفاوت، و گزارش‌های سنتز و تحلیل زمانی با حاشیهٔ مثبت، در مجموع درستی عملکرد و پیاده‌سازی‌پذیری طرح را تأیید می‌کنند.
\end{document}
